% --------------------------------- Preamble -----------------------------------
\documentclass[11pt, letterpaper, twocolumn]{article}

\usepackage[T1]{fontenc} % use 8-bit fonts
\usepackage[margin=0.75in, columnsep=0.2in, % 3.4 in columns
    headsep=0.28in, footskip=0.42in]{geometry} % paper sizes
\usepackage{amsmath, amsfonts, amssymb} % improved math
\usepackage{bm} % bold math (must be after amsmath)
\usepackage{upquote} % upright single quotes in verbatim
\usepackage{pgfplots} % includes tikz, xcolor, graphicx
\usepackage[colorlinks=true, allcolors=gray]{hyperref}
% \usepackage{float}

\raggedbottom % do not force stretching text to fill the page
\setlength{\emergencystretch}{\hsize} % avoid overfull lines
\overfullrule=1mm % black band to show overfull lines
\pdfsuppresswarningpagegroup=1 % remove pdf image warnings
\frenchspacing % consistent spacing between words
\newcommand{\ul}[1] % vectors
    {{}\mkern1mu\underline{\mkern-1mu#1\mkern-1mu}\mkern1mu}
\DeclareSymbolFont{euler}{U}{eur}{m}{n} % symbol font
\DeclareMathSymbol{\PI}{\mathord}{euler}{25} % up pi
\pgfplotsset{compat=newest} % use newest pgfplot version
\newcommand{\EE}[2]{\ensuremath{{{#1}\!\times\!10^{#2}}}}
% ------------------------------------------------------------------------------

\title{Nonlinear State Estimation of a Pendulum Using Vision Measurements \\ 
       \Large \textit{A comparative study of EKF, UKF, and Particle Filtering}}
\author{}
\date{} 

\begin{document}
\maketitle

\section*{Description}

This project addresses the problem of estimating the full state of a pendulum system 
using nonlinear dynamics and camera-based measurements. The system consists of a ball 
attached to a string, observed by a monocular camera that provides only image-space measurements.

The estimation problem is challenging due to:
\begin{itemize}
    \item nonlinear dynamics with multiple coupled forces
    \item nonlinear measurement model arising from perspective projection
    \item partial observability of the system (unknown hook position)
\end{itemize}

To study this problem, three Bayesian estimation methods are implemented and compared:
\begin{itemize}
    \item Extended Kalman Filter (EKF)
    \item Sigma-point Kalman Filter (UKF)
    \item Particle Filter (PF)
\end{itemize}

This project explores how different estimation strategies perform on the same nonlinear system,
highlighting trade-offs between accuracy, robustness, and computational complexity.

\section*{System Model}

The state vector and its dynamics are defined as
\[\begin{aligned}
\ul{x} &= \begin{bmatrix}
    \ul{p} & \ul{v} & \ul{h}
    \end{bmatrix}^{\top} \\
\dot{\ul{x}} &= \begin{bmatrix}
    \mathbf{0}_3 & \mathbf{I}_3 & \mathbf{0}_3 \\
    \mathbf{0}_3 & \mathbf{0}_3 & \mathbf{0}_3 \\
    \mathbf{0}_3 & \mathbf{0}_3 & \mathbf{0}_3
    \end{bmatrix}\ul{x} + 
    \begin{bmatrix}
        \ul{0}_3 \\
        \ul{a} \\
        \ul{0}_3
    \end{bmatrix}
\end{aligned}
\]
where $\ul{p}$ is the position of the ball, $\ul{v}$ is its velocity, and $\ul{h}$ 
is the position of the hook.

The dynamics include gravitational acceleration, string tension effects, and aerodynamic drag:
\[
\ul{a} = \ul{a}_g + \ul{a}_c + \ul{a}_d + \ul{g}
\]
Where we discuss each acceleration component below.

\subsection*{Counter-gravitational Acceleration}
We are assuming a non-elastic string, so any force gravity puts on the ball
will be countered by the string. The position of the ball, $\ul{p}$, relative
to the hook, $\ul{h}$, is
\[
\ul{s} = \ul{p} - \ul{h}
\]
Its unit vector is $\ul{u}_s = \ul{s}/s$, where $s = ||\ul{s}||$. So, the 
tension acceleration (counter-gravitational) on the string due to gravity is
\[
\ul{a}_g = -\ul{u}_s(\ul{u}_s \cdot \ul{g})
\]

\subsection*{Centripital Acceleration}
Another tension on the string is due to the velocity, $\ul{v}$, of the ball.
This is the centripetal acceleration.
\[
\ul{a}_c = -\ul{u}_s\frac{v^2}{s}
\]
where $v = ||\ul{v}||$.

\subsection*{Drag Acceleration}
As the ball moves through the air, there is a drag which opposes its motion.
This acceleration is defined in the opposite direction to the velocity and its
magnitude is the velocity squared scaled by the drag coefficient, $k_d$
\[
\ul{a}_d = -\ul{v}(k_dv)
\]

\subsection*{Discrete Space Dynamics}
In discrete space, the dynamics equations are
\[\begin{aligned}
\ul{x}_{k+1} &= \ul{\phi}(\ul{x}_k) \\
\ul{\phi}(\ul{x}) &= \mathbf{I}_9\ul{x}+T \begin{bmatrix}
    \mathbf{0}_3 & \mathbf{I}_3 & \mathbf{0}_3 \\
    \mathbf{0}_3 & \mathbf{0}_3 & \mathbf{0}_3 \\
    \mathbf{0}_3 & \mathbf{0}_3 & \mathbf{0}_3
    \end{bmatrix}\ul{x}+
    \begin{bmatrix}
        \frac{1}{2}T^2\ul{a} \\
        T\ul{a} \\
        \ul{0}_3
    \end{bmatrix}
\end{aligned}
\]

\subsection*{Dynamics Jacobians}
The Extended Kalman Filter requires the use of Jacobians to approxiamte how 
uncertainty in the state propagates through nonlinear dynamics by linearizing 
these dyamics. Below we will derive these Jacobians.

The Jacobian of the counter-gravitational acceleration is composed of
\[\begin{aligned}
    \frac{\partial\ul{a}_g}{\partial p_x} &= -2\frac{s_x}{s^2}\ul{a}_g-\frac{1}{s^2}
    \begin{bmatrix}
        2g_xs_x+g_ys_y+g_zs_z \\
        g_xs_y \\
        g_xs_z
    \end{bmatrix} \\
    \frac{\partial\ul{a}_g}{\partial p_y} &= -2\frac{s_y}{s^2}\ul{a}_g-\frac{1}{s^2}
    \begin{bmatrix}
        g_ys_x \\
        g_xs_x+2g_ys_y+g_zs_z \\
        g_ys_z
    \end{bmatrix} \\
    \frac{\partial\ul{a}_g}{\partial p_z} &= -2\frac{s_z}{s^2}\ul{a}_g-\frac{1}{s^2}
    \begin{bmatrix}
        g_zs_x \\
        g_zs_y \\
        g_xs_x+g_ys_y+2g_zs_z
    \end{bmatrix}
\end{aligned}
\]
The Jacobian of the centripetal acceleration is composed of
\[\begin{aligned}
    \frac{\partial\ul{a}_c}{\partial p_x} &= -\frac{v^2}{s^2}\left(
        \begin{bmatrix}
            1 \\
            0 \\
            0
        \end{bmatrix}-2\frac{s_x}{s^2}\ul{s}
    \right) &\quad
    \frac{\partial\ul{a}_c}{\partial v_x} &= -2\frac{\ul{s}}{s^2}v_x \\
    \frac{\partial\ul{a}_c}{\partial p_y} &= -\frac{v^2}{s^2}\left(
        \begin{bmatrix}
            0 \\
            1 \\
            0
        \end{bmatrix}-2\frac{s_y}{s^2}\ul{s}
    \right) &\quad
    \frac{\partial\ul{a}_c}{\partial v_y} &= -2\frac{\ul{s}}{s^2}v_y \\
    \frac{\partial\ul{a}_c}{\partial p_z} &= -\frac{v^2}{s^2}\left(
        \begin{bmatrix}
            0 \\
            0 \\
            1
        \end{bmatrix}-2\frac{s_z}{s^2}\ul{s}
    \right) &\quad
    \frac{\partial\ul{a}_c}{\partial v_z} &= -2\frac{\ul{s}}{s^2}v_z
    \end{aligned}
\]
The Jacobian of the drag acceleration is composed of
\[
\frac{\partial\ul{a}_d}{\partial\ul{v}} = -\frac{k_d}{v}(v^2\mathbf{I}_3+\ul{v}\ul{v}^{\top})
\]
Notice that these Jacobians have only been with respect to the position
and velocity of the ball. The Jacobians with respect to the hook position
are simply the negative of the Jacobians with respect to the ball position.

Putting all these together, the $\Phi$ matrix is
\[
\Phi = \mathbf{I}_9+T\begin{bmatrix}
    \mathbf{0}_3 & \mathbf{I}_3 & \mathbf{0}_3 \\
    \mathbf{0}_3 & \mathbf{0}_3 & \mathbf{0}_3 \\
    \mathbf{0}_3 & \mathbf{0}_3 & \mathbf{0}_3
    \end{bmatrix}+\begin{bmatrix}
        \frac{1}{2}T^2\frac{\partial\ul{a}}{\partial\ul{x}} \\
        T\frac{\partial\ul{a}}{\partial\ul{x}} \\
        \mathbf{0}_{3\times9}
    \end{bmatrix}
\]
The $\Phi$ matrix is the discretized dynamics Jacobian. The ${\partial\ul{a}/}{\partial\ul{x}}$
is a 3$\times$9 matrix. The first three columns are with respect to the ball position,
the next three with respect to the ball velocity, and the last three with respect
to the hook position. Put columns vectors like ${\partial\ul{a}_g}/{\partial p_x}$
into the first column, ${\partial\ul{a}_c}/{\partial v_x}$ into the fourth column,
and so on. Vectors like ${\partial\ul{a}_g}/{\partial p_x}$ and
${\partial\ul{a}_c}/{\partial p_x}$ would add.

\section*{Measurement Model}
The camera measurement vector is
\[
\ul{z} = \begin{bmatrix}
        r_p & u & v
    \end{bmatrix}^{\top}
\]

Measurements are obtained from a monocular camera. The sensor provides:
\begin{itemize}
    \item pixel coordinates $(u, v)$
    \item apparent radius $r_p$
\end{itemize}

These measurements are related to the 3D position through a nonlinear projection model:
\[ \begin{aligned}
r_p &= \kappa \frac{r_b}{p_z} \\
u &= c_x + \kappa \frac{p_x}{p_z} \\
v &= c_y + \kappa \frac{p_y}{p_z}
\end{aligned}
\]

Where $r_b$ is the true radius of the ball, $\kappa$ is a scaling factor, $c_x$ 
is the horizontal center of the camera image in pixels, and $c_y$ is the vertical 
center of the image in pixels. The measurement model introduces strong nonlinearities 
due to division by depth ($p_z$).

\subsection*{Measurement Model Jacobians}
Like the dynamics, the measurement model requires the use of Jacobians in
predicting how uncertainty in the state affects the predicted measurement.
Below we will derive the measurement Jacobian.

The measurement Jacobian is composed of
\[
\frac{\partial\ul{h}}{\partial\ul{p}} = \begin{bmatrix}
    0 & 0 & -\kappa\frac{r_b}{p_z^2} \\
    \frac{\kappa}{p_z} & 0 & -\kappa\frac{p_x}{p_z^2} \\
    0 & \frac{\kappa}{p_z} & -\kappa\frac{p_y}{p_z^2}
\end{bmatrix}
\]

\section*{Estimation Methods}

Three Bayesian estimation approaches are applied to the same system.

\subsection*{Extended Kalman Filter (EKF)}

The EKF linearizes the dynamics and measurement models using analytical Jacobians. This
provides computational efficiency but introduces approximation error in highly nonlinear regions.

\subsection*{Sigma-Point Kalman Filter (UKF)}

The sigma-point Kalman filter approximates nonlinear transformations using deterministically
selected sigma points. This avoids explicit linearization and more accurately captures uncertainty propagation.

\subsection*{Particle Filter (PF)}

The particle filter represents the state distribution using weighted samples. This allows it to 
capture non-Gaussian and multimodal distributions but requires significantly more computation.

\section*{Initialization and Implementation}

Initial position estimates are obtained by inverting the measurement model. Velocity 
is estimated using finite differences between early measurements.

The hook position is initialized based on prior geometric assumptions. Covariance matrices
and noise parameters are tuned empirically to achieve stable and accurate estimation.

Each filter is implemented using the same system and measurement models to enable direct comparison.

\section*{Results}

The performance of each estimator is evaluated using trajectory reconstruction and measurement consistency.

\subsection*{Trajectories}

Figures~\ref{fig:ekftraj}--\ref{fig:pftraj} show the estimated trajectories for each method.

\begin{figure}[h]
    \centering
    \includegraphics[width=\linewidth]{../figures/fig_path_ekf.pdf}
    \caption{Estimated trajectory for EKF}
    \label{fig:ekftraj}
\end{figure}

\begin{figure}[h]
    \centering
    \includegraphics[width=\linewidth]{../figures/fig_path_ukf.pdf}
    \caption{Estimated trajectory for UKF}
    \label{fig:ukftraj}
\end{figure}

\begin{figure}[h]
    \centering
    \includegraphics[width=\linewidth]{../figures/fig_path_particle_filter.pdf}
    \caption{Estimated trajectory for PF}
    \label{fig:pftraj}
\end{figure}

\clearpage

Figures~\ref{fig:ekfbpos}--\ref{fig:pfbpos} show the estimated ball positions for each method.

\begin{figure}[h]
    \centering
    \includegraphics[width=0.72\linewidth]{../figures/fig_ball_position_ekf.pdf}
    \caption{Estimated ball position for EKF}
    \label{fig:ekfbpos}
\end{figure}

\begin{figure}[h]
    \centering
    \includegraphics[width=0.72\linewidth]{../figures/fig_ball_position_ukf.pdf}
    \caption{Estimated ball position for UKF}
    \label{fig:ukfbpos}
\end{figure}

\begin{figure}[h]
    \centering
    \includegraphics[width=0.72\linewidth]{../figures/fig_ball_position_particle_filter.pdf}
    \caption{Estimated ball position for PF}
    \label{fig:pfbpos}
\end{figure}

\pagebreak

Figures~\ref{fig:ekfbvel}--\ref{fig:pfbvel} show the estimated ball velocities for each method.

\begin{figure}[h]
    \centering
    \includegraphics[width=0.72\linewidth]{../figures/fig_ball_velocity_ekf.pdf}
    \caption{Estimated ball velocity for EKF}
    \label{fig:ekfbvel}
\end{figure}

\begin{figure}[h]
    \centering
    \includegraphics[width=0.72\linewidth]{../figures/fig_ball_velocity_ukf.pdf}
    \caption{Estimated ball velocity for UKF}
    \label{fig:ukfbvel}
\end{figure}

\begin{figure}[h]
    \centering
    \includegraphics[width=0.72\linewidth]{../figures/fig_ball_velocity_particle_filter.pdf}
    \caption{Estimated ball position for PF}
    \label{fig:pfbvel}
\end{figure}

\clearpage

Figures~\ref{fig:ekfhpos}--\ref{fig:pfhpos} show the estimated hook positions for each method.

\begin{figure}[h]
    \centering
    \includegraphics[width=0.72\linewidth]{../figures/fig_hook_position_ekf.pdf}
    \caption{Estimated hook position for EKF}
    \label{fig:ekfhpos}
\end{figure}

\begin{figure}[h]
    \centering
    \includegraphics[width=0.72\linewidth]{../figures/fig_hook_position_ukf.pdf}
    \caption{Estimated hook position for UKF}
    \label{fig:ukfhpos}
\end{figure}

\begin{figure}[h]
    \centering
    \includegraphics[width=0.72\linewidth]{../figures/fig_hook_position_particle_filter.pdf}
    \caption{Estimated hook position for PF}
    \label{fig:pfhpos}
\end{figure}

\pagebreak

\subsection*{IOU}

Performance is also evaluated using the IOU metric between predicted and observed 
image measurements. This function calculates the intersect-over-union (IOU) 
of the expected pixel cirlce and the measured pixel circle. If there is  a perfect 
overlap, the  IOU value will be 1. Anything else will return less than 1 (minimum 0).

Table~\ref{fig:iou} shows the average IOU over the duration of the simulation
for each filtering method.

\begin{table}[h]
    \centering
    \begin{tabular}{|c|c|}
        \hline
        Method & IOU \\ \hline
        EKF & 0.99733 \\ \hline
        UKF & 0.99715 \\ \hline
        PF & 0.97059 \\ \hline
    \end{tabular}
\caption{Intersection-over-union (IOU) values}
\label{fig:iou}
\end{table}

\section*{Discussion}

The results highlight clear differences in performance across the three estimation methods.

The EKF achieved the highest IOU, indicating that for this problem the linearization-based 
approach was sufficiently accurate. This suggests that, despite the nonlinear measurement model,
the system operates in a manner where local linear approximations remain valid over each update step.

The UKF performed very similarly to the EKF, with only a slightly lower IOU. This indicates
that the additional complexity of sigma-point propagation did not provide a significant advantage
in this specific scenario. However, the UKF still offers improved robustness to stronger nonlinearities
and does not require analytical Jacobians, which makes it a more flexible approach for more complex systems.

The particle filter produced a lower IOU compared to both Kalman filter approaches, though its 
performance remained strong overall. Achieving this level of accuracy required approximately 500 
particles, highlighting the trade-off between computational cost and estimation quality. While the 
particle filter is capable of handling highly nonlinear and non-Gaussian systems, its increased 
computational burden makes it less efficient for problems where Gaussian assumptions are reasonable.

An important observation is that all three filters were able to infer the full 9-dimensional state 
from indirect camera measurements. This demonstrates that the combination of system dynamics and 
measurement structure provides sufficient information to reconstruct the underlying state, even when 
certain variables (such as depth and hook position) are not directly measured.

Overall, the results suggest that for this problem, the system is sufficiently well-behaved that 
Gaussian-based filters (EKF and UKF) are both accurate and efficient, while the particle filter 
provides additional robustness at a higher computational cost.

\section*{Conclusion}

This project investigated nonlinear state estimation for a pendulum system using vision-based measurements. 
Three estimation methods were implemented and compared: an Extended Kalman Filter, a Sigma-Point Kalman 
Filter (UKF), and a Particle Filter.

The results show that all three methods are capable of accurately reconstructing the system state, 
despite the nonlinear dynamics and indirect measurements. The EKF achieved the highest performance in 
terms of IOU, with the UKF performing nearly as well. The particle filter also produced accurate estimates 
but required significantly greater computational effort to achieve comparable performance.

These results highlight an important principle in state estimation: the choice of estimator should be guided 
by the level of system nonlinearity and the required computational efficiency. For systems where nonlinearities 
are moderate and Gaussian assumptions are reasonable, Kalman filter-based methods provide an effective and 
efficient solution. More general methods such as particle filters become advantageous when these assumptions break down.

This project demonstrates the full estimation pipeline for a nonlinear, partially observed system, 
including modeling, measurement inversion, filter design, and performance evaluation. It also illustrates 
how different estimation approaches behave when applied to the same problem, providing insight into their 
relative strengths and limitations.

\end{document}